% =====================================================================
% Supplementary Information: Cross-lingual sentiment validation
% =====================================================================
%
% Compile-time dependencies (add to preamble if not already present):
%   \usepackage{booktabs}
%   \usepackage{siunitx}
%   \usepackage{amsmath}
%   \usepackage{graphicx}
%
% Drop this subsection into the SI \section{Supplementary Information}.

\subsection{Cross-lingual sentiment validation and correction}
\label{si:cross-lingual}

\subsubsection{Motivation: why two different scoring models}

The Weibo (China) and Twitter (USA) opinion time series in the main text
were originally produced by two different large-language-model analyzers:
Weibo posts (Chinese) were scored by \textbf{Kimi-K2} with a
Chinese-language prompt, while tweets (English) were scored by
\textbf{GPT-5} with an English-language prompt. This split was driven by a
data-access constraint rather than a methodological preference: the Weibo
collection pipeline runs on a server inside mainland China, where direct
calls to the OpenAI API are not permitted. Routing every Weibo post out of
that environment for re-scoring was infeasible at the scale of our
corpus. Conversely, the Twitter pipeline lives outside China and can
freely call GPT-5. At the time of the original analysis, Kimi-K2 and
GPT-5 were within roughly one percentage point of each other on
publicly reported reasoning and instruction-following benchmarks, so the
choice was made to use whichever model was natively available on each
side.

A reviewer pointed out that even when two models perform comparably on
generic benchmarks, they may apply systematically different sentiment
scales to the same text. We therefore conducted a cross-lingual,
cross-model validation experiment to (i) estimate the residual offset
between the two analyzers when both score the \emph{same content}, and
(ii) apply that estimate as an additive correction to the Weibo line in
the comparison figures.

\subsubsection{Experimental design}

We instantiated a symmetric two-arm experiment, one per platform:

\paragraph{Sampling.} From each platform we drew a stratified random
sample of $n = 1000$ posts (seed fixed; proportional stratification on
the original opinion class). Weibo posts retained their original
Chinese text; tweets retained their original English text.

\paragraph{Machine translation.} Each sample was translated once
through the DeepL API: Weibo posts zh$\rightarrow$en, tweets
en$\rightarrow$zh. Translations were cached locally so that all
subsequent rounds operated on identical translated text and any
variance came from the language model alone, not from sampling or
translation noise.

\paragraph{Cross-lingual re-scoring.} The translated text on each side
was then re-scored by the \emph{other} analyzer:
translated Weibo (now in English) was re-scored by GPT-5 with the
English prompt used for tweets; translated tweets (now in Chinese) were
re-scored by Kimi-K2 with the Chinese prompt used for Weibos. To
distinguish model sampling noise from systematic bias, the re-scoring
was repeated $R = 5$ independent rounds with identical inputs.

\paragraph{Per-round signed difference.} For each round $r \in
\{1,\dots,5\}$ and each platform $p \in \{\text{weibo},
\text{tweet}\}$, we restricted to posts that returned a valid integer
score on both the original and the cross-lingual side, then computed
the mean signed difference
\begin{equation*}
\Delta^{(p)}_r \;=\; \frac{1}{|\mathcal{V}^{(p)}_r|} \sum_{i \in
\mathcal{V}^{(p)}_r} \bigl(\,o^{\text{cross}}_{i,r} -
o^{\text{orig}}_i\bigr),
\end{equation*}
where $o^{\text{orig}}_i$ is the post's original score and
$o^{\text{cross}}_{i,r}$ is its round-$r$ cross-lingual re-score, both
on the five-point opinion scale $\{-2,-1,0,1,2\}$.

\subsubsection{Correction factor}

The two arms provide independent estimates of the same Kimi-Chinese
vs.\ GPT-English scoring offset, but they are measured on different
corpora. Because the quantity we ultimately correct is the
\emph{Weibo} time series, we use the Weibo-arm estimate directly: it is
the offset measured on the exact content domain that the correction is
applied to. It is also the larger (more conservative) of the two arms
in absolute value, so this choice does not understate the residual
bias.

Concretely, we define the round-averaged Weibo-arm offset
\begin{equation}
\widehat{\Delta}^{(\text{weibo})} \;=\; \frac{1}{R} \sum_{r=1}^{R}
\Delta^{(\text{weibo})}_r,
\label{eq:correction}
\end{equation}
which directly estimates the mean of $\text{GPT-en}(x) -
\text{Kimi-zh}(x)$ across Weibo posts $x$. We then add this offset
(with its sign) to every value of the original Weibo opinion time
series before plotting:
\begin{equation}
y^{\text{corrected}}_{\text{weibo}}(t) \;=\; y^{\text{orig}}_{\text{weibo}}(t)
+ \widehat{\Delta}^{(\text{weibo})}.
\label{eq:apply}
\end{equation}
Interpreted operationally,
Equation~\eqref{eq:apply} translates the Kimi-Chinese-scored Weibo line
onto the GPT-English scale used for the Twitter line, making the two
series directly comparable. The correction is a vertical shift only:
the rolling-mean smoothing in the comparison figure is applied to the
raw daily values first, and the correction is applied as a final
additive offset, so temporal dynamics are unchanged.

The Twitter-arm estimate
$\widehat{\Delta}^{(\text{tweet})} = +0.0464$ is reported in
Table~\ref{tab:cross-lingual-rounds} as an independent cross-check:
both arms agree in direction (Kimi-Chinese assigns slightly more
positive opinion scores than GPT-English on the same translated
content), confirming that the offset is genuine rather than an
artifact of one arm's content domain.

\subsubsection{Results}

Table~\ref{tab:cross-lingual-rounds} reports, for each of the five
rounds, the number of posts that yielded a valid score on both sides,
and the per-round signed differences $\Delta^{(\text{weibo})}_r$ and
$\Delta^{(\text{tweet})}_r$.

\begin{table}[h]
\centering
\caption{Cross-lingual signed sentiment differences per round.
$\Delta^{(\text{weibo})}_r$: GPT-English re-scores of translated Weibo
posts minus their original Kimi-Chinese scores. $\Delta^{(\text{tweet})}_r$:
Kimi-Chinese re-scores of translated tweets minus their original
GPT-English scores. $n^{(p)}_r$ is the count of valid
original--cross-lingual score pairs (out of $1000$ sampled posts per
platform) in round $r$. The applied correction is the round-averaged
Weibo-arm value, shown in bold.}
\label{tab:cross-lingual-rounds}
\begin{tabular}{c S[table-format=4.0] S[table-format=4.0]
                  S[table-format=+1.4] S[table-format=+1.4]}
\toprule
Round &
{$n^{(\text{weibo})}_r$} &
{$n^{(\text{tweet})}_r$} &
{$\Delta^{(\text{weibo})}_r$} &
{$\Delta^{(\text{tweet})}_r$} \\
\midrule
1 & 848 & 753 & -0.0814 & +0.0438 \\
2 & 857 & 767 & -0.0782 & +0.0495 \\
3 & 849 & 755 & -0.0789 & +0.0477 \\
4 & 856 & 752 & -0.0537 & +0.0439 \\
5 & 840 & 763 & -0.0690 & +0.0472 \\
\midrule
\textbf{Mean} & \textbf{850.0} & \textbf{758.0} &
\textbf{-0.0722} & {+0.0464} \\
SD (across rounds) & {6.9} & {6.6} & {0.0114} & {0.0025} \\
\bottomrule
\end{tabular}
\end{table}

Both arms produce small differences in absolute value ($|\Delta| <
0.10$ on the five-point scale) that are remarkably stable across rounds
(SDs of $0.011$ and $0.003$ respectively), indicating that the
disagreement between the two analyzers is a systematic offset rather
than per-call sampling noise. The applied correction is
$\widehat{\Delta}^{(\text{weibo})} = -0.0722$ on the $[-2,+2]$ opinion
scale.

\subsubsection{Producing the corrected figure}

The comparison plot script (\texttt{plot.py}) exposes the correction as
a CLI flag. The uncorrected and corrected figures are produced as
follows:

\begin{verbatim}
# Uncorrected (Weibo line on Kimi-Chinese scale)
python plot.py --apply_correction False

# Corrected (Weibo line shifted onto GPT-English scale)
python plot.py --apply_correction True --correction_factor -0.0722
\end{verbatim}

The filename of the saved PDF encodes the correction state
(\texttt{\_uncorrected} or \texttt{\_corrected\_weibo-0.0722}), so the
two variants can be regenerated and compared without overwriting each
other.
